\chapter{Methodology}
\label{ch:method}

This chapter specifies how A.L.F.R.E.D.\ is evaluated, so that every number in
Chapter~\ref{ch:eval} can be read against a fixed and stated procedure. Two
methodological commitments do most of the work and are described first: scoring
only the execution stage in isolation (\emph{binding-only},
Section~\ref{sec:m-binding}) and scoring success by the resulting device state
rather than the command string (the \emph{end-state oracle},
Section~\ref{sec:m-oracle}). The benchmarks, the safeguards against runaway
agents, and the integrity rules follow.

\section{Models and infrastructure}
\label{sec:m-models}

Several models appear in the experiments, all run locally. The project began with
\texttt{ministral-3-3b} as the reflexer and later adopted \texttt{qwen-3.5-2b}
after the model sweep of Section~\ref{sec:p0}; both are used as command-generating
reflexers, and \texttt{qwen-3.5-2b} additionally serves as the small-model
\textbf{thinker} baseline. The large \textbf{thinker} is \texttt{qwen3.6-35b}, a
mixture-of-experts reasoning model (35B total, 3B active) used as the capability
ceiling for free-form binding and as a local distillation teacher. The cloud
distillation teacher is Claude Sonnet~4.6. The Qwen3.5 family at 0.8B and 4B also
appears in the reflexer sweep.

\paragraph{Serving stack and quantization.} Local models are served with
\texttt{llama-server} from \texttt{llama.cpp}~\cite{llamacpp} over its
OpenAI-compatible HTTP endpoint, and the agent loop is driven by
\texttt{pi}~\cite{pi-agent}. \textbf{All models use 4-bit (Q4) GGUF
quantization}, so that every reported number reflects a model that fits a
phone-class memory budget. The two executor models (\texttt{qwen-3.5-2b} and
\texttt{ministral-3-3b}) are served with \textbf{identical decoding parameters};
the qwen-3.5-2b invocation is reproduced verbatim in Listing~\ref{lst:serve} and
ministral-3-3b uses the same flags. The decisive settings are a low sampling
temperature ($0.2$) for near-deterministic command generation and, crucially,
\textbf{reasoning disabled} (\texttt{enable\_thinking=false})---the reflexer is an
executor, not a reasoner, so chain-of-thought is switched off. The 35B is run in
its native reasoning mode (thinking enabled), since it is evaluated precisely as
the free-form reasoning ceiling.

\begin{lstlisting}[caption={Reflexer serving configuration (\texttt{qwen-3.5-2b});
\texttt{ministral-3-3b} uses the same parameters. Q4 quantization, temperature
$0.2$, reasoning disabled.},label={lst:serve},basicstyle=\ttfamily\scriptsize]
llama-server Qwen3.5-2B-UD-Q4_K_XL.gguf --alias qwen-3.5-2b
  --host 0.0.0.0 --port 8005 --ctx-size 4072 --n-gpu-layers 999 --jinja
  --threads 16 --temp 0.2 --top-p 0.95 --top-k 20 --min-p 0.00
  --presence_penalty 0.5 --cache-type-k q8_0 --cache-type-v q8_0
  --metrics --slots --flash-attn on --cache-ram 0
  --chat-template-kwargs '{"enable_thinking":false}'
\end{lstlisting}

Runs are executed \emph{sequentially}: concurrent traffic to the same endpoint was
observed to corrupt token accounting and depress accuracy
(Section~\ref{sec:m-integrity}). Model identities, ports, and per-run parameters
are recorded with every report under \texttt{benchmark/reports/}.

\section{Binding-only evaluation}
\label{sec:m-binding}

End-to-end accuracy is the product of two stages: \emph{routing} (does the system
pick the right pattern or skill?) and \emph{binding} (given the right target,
does it emit the right command?). To measure how much model the execution stage
needs, routing must be removed as a confound. The binding-only protocol does
this by assuming perfect routing: the reflexer is handed the gold pattern for
each task, and the thinker is given the documentation for the correct skill.
Only binding is then under test. This isolation is what licenses the comparisons
in Sections~\ref{sec:p1}--\ref{sec:p2}; the cost of treating routing separately
is revisited honestly in Section~\ref{sec:routing}, where it is identified as the
real end-to-end bottleneck.

For the binding experiments, success is scored by \emph{token coverage}: a task
carries a set of gold tokens that must appear (case-insensitively) in the union
of the commands the model emitted. Both tracks read the same task files and use
the same checker---the reflexer harness
(\texttt{alfred/eval/binding\_only\_eval.py}) and the thinker harness
(\texttt{src/run-binding-compare.ts}) implement identical coverage semantics---so
that the only difference between conditions is the model and whether it received
the pattern. Standardizing the checker across both tracks was a deliberate change
to make the headline comparison fair: the reflexer is not scored by an easier
rule than the thinker.

\section{The end-state oracle}
\label{sec:m-oracle}

Token coverage is adequate for binding, where commands are short and singular,
but it is the wrong instrument for the settings benchmarks, where a request can
be satisfied by more than one command and a syntactically valid command can still
produce the wrong outcome. The settings benchmarks (SET-1 and SET-2) therefore
use a different, stricter criterion: an \textbf{end-state oracle}. Each task runs
against a mutable state backend; after the agent finishes, the oracle compares
the final state to a gold end-state and passes the task only if two conditions
hold: the requested keys reached their target values, \emph{and} no other key was
changed (no collateral edits).

This two-part definition catches two failures that command matching misses. The
first is ``right command, wrong outcome''---a command that looks correct but
leaves the device in the wrong state. The second is over-action---an agent that
achieves the goal but also changes settings it was never asked to touch, which a
command-matching scorer would happily pass. The no-collateral clause is not
hypothetical: in the distillation experiments an over-verbose pattern caused the
student to toggle extra radios, and only the end-state oracle detected it
(Section~\ref{sec:p5}). The oracle is implemented once per benchmark
(\texttt{SET-1/oracle.py}, \texttt{SET-2/oracle.py}) as the single source of
truth for success, following the contract-first principle of
Section~\ref{sec:principles}: the oracle and the task contract were written
before the runner.

\section{Sandbox isolation and the kill policy}
\label{sec:m-sandbox}

Every settings task runs in a fresh sandbox. The backend state is materialized
into a temporary file seeded with the task's initial state, the agent acts only
on that copy, and the sandbox is destroyed afterward. No task can see another
task's changes, and no run touches real device state, the network, or any
persistent user data. If an agent corrupts the state file (for example by writing
malformed JSON), the oracle reports the corruption rather than crashing, so the
failure is recorded as a failure rather than lost.

Agents that reason in a loop must be prevented from running forever. Two
mechanisms bound every run. A \emph{step ceiling} caps the number of commands an
agent may issue: for atom binding the ceiling is twice the expected number of
moves (an atom should need one command, so the cap is two), and for the
exploratory settings runs the ceiling is larger to permit genuine discovery. A
\emph{wall-clock backstop} aborts a run that exceeds a time budget regardless of
step count, which matters for the slow 35B agent. When a ceiling is hit the agent
is aborted and scored on whatever state it has produced, so a non-terminating
agent counts as a (possibly partial) failure rather than stalling the benchmark.
Both bounds, and the commands each agent actually issued, are recorded in the run
report for inspection.

\section{Benchmarks}
\label{sec:m-benchmarks}

Three benchmarks are used, chosen to bracket the difficulty of on-device action.

\paragraph{Expansion skills (binding).} Seven execution skills---clock, media,
focus, device settings, call, app, and browser---with thirty-six atoms and a set
of cross-skill composites, each backed by a mock CLI that emits a structured
result. The task set provides one to two queries per skill (the
\emph{balanced-10}) and a larger per-skill set (seventy-two atom tasks), with gold
tokens for coverage scoring. This benchmark isolates binding across a realistic
skill surface and is the basis for Sections~\ref{sec:p1}--\ref{sec:p2}.

\paragraph{SET-1 (clean settings API).} A friendly settings interface
(\texttt{settings set connectivity.wifi off}) with transparent intent-to-path
mapping and validation that rejects bad values. Twenty canonical tasks plus ten
harder multi-step tasks, scored by the end-state oracle. SET-1 isolates the value
of a deep tool on an easy surface.

\paragraph{SET-2 (realistic Android API).} An \texttt{adb}-style interface
(\texttt{settings put <ns> <key> <value>}) over thirty-two settings in three
namespaces (\texttt{system}, \texttt{secure}, \texttt{global}), with non-obvious
encodings---brightness on 0--255, timeouts in milliseconds, integer enums for Do
Not Disturb, ringer, and location, and a string-float font scale---and a
permissive \texttt{put} that silently accepts wrong inputs, exactly as the real
Android tool behaves. Forty-eight tasks, scored by the end-state oracle. SET-2 is
the realistic stressor and the setting for the limit-of-scale and distillation
results (Sections~\ref{sec:p4}--\ref{sec:p5}).

\section{Integrity protocol}
\label{sec:m-integrity}

Because the central claims rest on a handful of comparisons, the evaluation
follows explicit integrity rules, stated here so they can be checked.

\begin{enumerate}[leftmargin=2em]
  \item \textbf{No fabricated numbers.} Every figure and table cites a saved run
    under \texttt{benchmark/reports/}; no result is estimated, rounded for
    effect, or hand-edited. Where two runs of the same condition disagree, both
    are reported rather than reconciled silently.
  \item \textbf{No benchmark leakage into distilled artifacts.} In the
    distillation experiments the teacher is given only the skill documentation
    and the API schema, never any task or its gold outcome. A pattern, once
    minted, is \emph{frozen}: it is never edited after observing which tasks it
    failed, because editing it would leak the benchmark into the artifact under
    test. The frozen local-teacher patterns are committed as evidence
    (\texttt{SET-2/patterns\_distilled\_pi.json}).
  \item \textbf{Isolation.} Runs never touch real user state, never call the
    device's execution-logging path, and never use the network; each settings
    task runs in a destroyed-after-use sandbox (Section~\ref{sec:m-sandbox}).
  \item \textbf{Sequential execution.} Concurrent runs against the same model
    endpoint were found to corrupt token accounting and depress accuracy, so the
    reported runs are executed one at a time.
\end{enumerate}

These rules are the methodological backbone of the thesis: they are what allow a
small number of comparisons to carry a strong claim. The next chapter applies the
procedure defined here and reports what it found.
